% krylov_trio_close.tex --- companion technical note to
% paper/krylov_diameter/krylov_diameter.tex
%
% Trio closure attempt for Krylov-Diameter Thm 4 via
% Vardian (2602.02675) + Heller-Papalini-Schuhmann (2412.17785, PRL 135, 151602)
% + Leutheusser-Liu (2508.00056). VERDICT: chain fails on three independent obstructions.
%
% Author: Kevin Remondiere (drafted by Opus 4.7 1M-context audit, 2026-05-02)
% Compile: pdflatex krylov_trio_close

\documentclass[11pt]{article}

\usepackage[utf8]{inputenc}
\usepackage[T1]{fontenc}
\usepackage[margin=1in]{geometry}
\usepackage{amsmath, amssymb, amsthm}
\usepackage{mathtools}
\usepackage{hyperref}
\usepackage{xcolor}
\usepackage{booktabs}
\hypersetup{colorlinks=true, linkcolor=blue, citecolor=blue, urlcolor=blue}

\newtheorem{obstruction}{Obstruction}
\newtheorem{remark}{Remark}

\newcommand{\R}{\mathbb{R}}
\newcommand{\HH}{\mathcal{H}}
\newcommand{\AA}{\mathcal{A}}
\newcommand{\NN}{\mathcal{N}}
\newcommand{\MM}{\mathcal{M}}
\newcommand{\OO}{\mathcal{O}}
\newcommand{\HL}{\mathrm{HL}}
\newcommand{\FRW}{\mathrm{FRW}}
\newcommand{\Rprop}{R_{\mathrm{proper}}}

\title{Trio Closure Attempt for Krylov--Diameter Theorem~4: \\
Vardian $\oplus$ Heller--Papalini--Schuhmann $\oplus$ Leutheusser--Liu --- Verdict and Obstructions}
\author{Kevin Remondi\`ere}
\date{2026-05-02}

\begin{document}
\maketitle

\begin{abstract}
We test whether the recent trio of published results --- Vardian
(arXiv:2602.02675, hep-th, 2026), Heller--Papalini--Schuhmann
(arXiv:2412.17785, \emph{Phys.\ Rev.\ Lett.}\ \textbf{135} (2025) 151602),
and Leutheusser--Liu (arXiv:2508.00056, hep-th, 2025) --- supplies
the three published links needed to upgrade Theorem~4
(Krylov--Diameter Correspondence) of the companion note
\texttt{paper/krylov\_diameter/krylov\_diameter.tex} from
``conditional on UOGH transfer'' to unconditional. We verify all three
papers exist on arXiv (opensearch \texttt{totalResults}=1 each), read
each in full, and conclude that the chain
\emph{fails} on three independent obstructions (T1, T2, T3) plus two
secondary ones (T5, T6, the latter inherited from the prior
Vardian--Hollands attempt). The Sonnet-frontier-search composite
claim ``modular Krylov spread complexity = log Jones-index of
entanglement-wedge subalgebra inclusion = holographic
complexity-volume'' is misattributed: HPS proves
\emph{Krylov-spread $=$ bulk volume} in DSSYK/sine-dilaton gravity, not
\emph{Krylov-spread $=$ log Jones-index}; only Leutheusser--Liu
supplies the second equality, and they do so by \emph{defining} a new
holographic index rather than identifying an existing
subfactor invariant. We recommend adding a defensive-citation
appendix to \texttt{krylov\_diameter.tex} listing T1--T6 with the
verbatim quotes below.
\end{abstract}

\section{Verbatim theorems with equation numbers}\label{sec:verbatim}

\paragraph{Vardian arXiv:2602.02675, eq.~(26).}
Verbatim (p.~4):
\begin{quote}
``$O\in\AA \to \{a^O_n,b^O_n\} \to (\mathcal{L}_S)_{mn} = [\text{tridiagonal Lanczos matrix}] \to P_O\,K|_{\AA}\,P_O : \text{spectrum of } K|_{\AA}\to \text{Area Operator}$.''
\end{quote}
with structural identity (eq.~27):
$\mathcal{L}_A = \bigoplus_\alpha S(\chi_\alpha,\MM_A)P_\alpha = K_A \otimes I_{\bar A} - K|_\AA$.
Operator modular Liouvillian (eqs.~33--34):
$L_S O := [K_A \otimes I_{\bar A}, O]$.

\paragraph{Heller--Papalini--Schuhmann arXiv:2412.17785 = \emph{Phys.~Rev.\ Lett.}\ \textbf{135} (2025) 151602, eq.~(13).}
Verbatim (p.~3):
\begin{quote}
``$2\,|\log q|\, C_K(t)_\beta = \langle \hat{L}\rangle = [-\partial_\Delta\langle e^{-\Delta\hat L}\rangle]_{\Delta=0}$.''
\end{quote}
Operator form (eq.~30): $2\,|\log q|\, \hat K = \hat L$. The state is
$|\psi(t)\rangle_\beta = Z_\beta^{-1/2} e^{-iT(t-i\beta/2)} |0\rangle$ (eq.~32) on the DSSYK
chord Hilbert space, dual to sine dilaton gravity (action eq.~2).
\textbf{Setting:} 2D dilaton gravity / finite-$N$ DSSYK. \emph{No mention} of Jones
index, log index, subfactor, Kosaki, type II$_\infty$, FRW, or Murray--von~Neumann
index anywhere in the 9-page letter (verified by full PDF read).

\paragraph{Leutheusser--Liu arXiv:2508.00056, eq.~(4).}
Verbatim:
\begin{quote}
``$(\MM:\NN) = \exp\!\big(C\cdot\mathrm{Vol}(\mathfrak{b})\big)$.''
\end{quote}
where $(\MM:\NN)$ is a \emph{new holographic index} (not Kosaki --- explicit disclaimer:
\emph{``Currently, it is not clear if (4) can be related to the Kosaki index.
Instead, we may view (4) as a definition of a new index using holography.''}),
$\mathfrak{b}$ is the maximal-volume slice of the bulk subregion dual to relative
commutant $\NN'\cap\MM$, and $C\propto 1/G_N$. \textbf{Setting:} AdS subregion-subalgebra
duality; type III$_1$ in large-$N$. \S{}V applies to dS observer algebras at the
North/South poles; FRW is not discussed.

\section{Composite chain attempt}\label{sec:chain}

On the type II$_\infty$ FRW comoving diamond
$\NN := \AA(D_O)_\FRW \rtimes_\sigma \R$ (\cite[Thm.~3.5]{Remondiere2026}):
\begin{itemize}
\item \textbf{Link A (Vardian).} Lanczos-extract $\Spec(K|_\AA)$ from boundary modular
Krylov data via eq.~(26).
\item \textbf{Link B (HPS).} Convert Krylov spread complexity to bulk geodesic length
via eq.~(13): $2|\log q|\, C_K = \langle\hat L\rangle$.
\item \textbf{Link C (LL).} Identify $\mathrm{Vol}(\mathfrak{b}) = (1/C)\log(\MM:\NN)$
via eq.~(4).
\end{itemize}
Composing B$\circ$C: $C_K \propto \mathrm{Vol} \propto \log(\MM:\NN)$. Combine with A
to identify the area operator on the boundary side as the central element of $K|_\AA$.

\section{Verdict: the chain does not close}\label{sec:verdict}

\paragraph{Status: chain DOES NOT CLOSE Krylov--Diameter Theorem~4.}
Three independent obstructions, each decisive:

\begin{obstruction}[HPS does not contain a log-index identity]\label{obs:T1}
HPS (arXiv:2412.17785, PRL 135, 151602) does not say
``Krylov-spread $=$ log Jones-index''. Verified by full read of pages 1--9: zero
occurrences of ``Jones'', ``subfactor'', ``Kosaki'', ``log index'', ``type
II$_\infty$ factor'', ``Murray--von~Neumann'', or ``FRW''. HPS proves a
\emph{complexity-volume} match in 2D (sine dilaton $\leftrightarrow$ DSSYK transfer
matrix), with the volume being the geodesic length $\hat L$ of the Kruskal extension
of the AdS$_2$ black hole (eq.~4). The Sonnet trio's middle link
``$=$ log Jones-index'' is misattributed --- HPS supplies ``Krylov $=$ volume'',
not ``Krylov $=$ log index''. Only Leutheusser--Liu supplies ``volume $=$ log
index''. So the Sonnet chain is really a TWO-link chain HPS $\circ$ LL, not a
triple, and Vardian's OAQEC area-operator construction is an independent piece
that does not compose with the HPS/LL line.
\end{obstruction}

\begin{obstruction}[Setting mismatch: incompatible algebra types]\label{obs:T2}
\begin{center}
\renewcommand{\arraystretch}{1.2}
\begin{tabular}{p{2.4cm}p{4.4cm}p{4.0cm}p{3.6cm}}
\toprule
Paper & Algebra & Setting & vs.\ KD target \\
\midrule
Vardian   & type-I code, OAQEC, $\bigoplus_\alpha L(\HH_A^\alpha)\otimes I$ & AdS/CFT QES + islands & II$_\infty$ FRW crossed product \\
HPS       & DSSYK chord Hilbert space (separable, transfer-matrix Lanczos) & 2D sine dilaton gravity & II$_\infty$ FRW (4D conf.\ flat) \\
LL        & type III$_1$ large-$N$, ``new holographic index'' (not Kosaki) & AdS subregion-subalgebra + dS poles & II$_\infty$ FRW with semifinite trace \\
\bottomrule
\end{tabular}
\end{center}
None of the three is on a type II$_\infty$ crossed product of an FRW
comoving-diamond algebra. Vardian fails by O1 of the prior
Vardian--Hollands run (OAQEC tensor decomposition fails in II$_\infty$).
HPS's chord Hilbert space is intrinsic to DSSYK and not equivalent to
$L^2(\NN,\mathrm{tr}_\NN)$ for any FRW crossed product; the holographic dictionary
$L = 2|\log q|n$ is specific to q-deformed JT/sine-dilaton. LL is the closest
(III$_1$ in large-$N$ matches the underlying III$_1$ algebra
$\AA(D_O)_\FRW$ from \cite[Thm.~3.5]{Remondiere2026}), but their ``index''
is \emph{defined} holographically via eq.~(4), not derived from a subfactor
calculation, and their AdS bulk subregion is not the FRW comoving diamond.
\end{obstruction}

\begin{obstruction}[LL ``index'' is a definition, making the chain circular on FRW]\label{obs:T3}
LL themselves disclaim that (4) is a theorem on Kosaki/Jones index. For an
AdS bulk subregion they have an \emph{independent} definition of
$\mathrm{Vol}(\mathfrak{b})$ (the maximal-volume slice in AdS) and they
\emph{postulate} the index by exponentiation. To use (4) as the right-hand link
of a chain that is supposed to \emph{prove} the Krylov--Diameter rate
$(1/C_k)\,\mathrm{d}C_k/\mathrm{d}t = 1/\Rprop$ on FRW, one would need either
(i) an independent, subfactor-theoretic computation of an index for the
inclusion $\AA(D_O')_\FRW \subset \AA(D_O)_\FRW$ of nested FRW comoving
diamonds, or (ii) an independent geometric meaning of $\mathrm{Vol}(\mathfrak{b})$
for the FRW diamond. Neither is established in LL or in the FRW companion
note. Plugging (4) in as a \emph{definition} makes the resulting
``Krylov $=$ volume'' identity tautological, not a theorem.
\end{obstruction}

\begin{obstruction}[Volume mismatch HPS$\,\nleftrightarrow\,$LL]\label{obs:T5}
HPS uses the volume of a Kruskal-extended AdS$_2$ BH slice (eq.~4); LL uses the
maximal-volume slice of an AdS$_d$ bulk subregion. Neither volume is
$\Rprop(\eta_c)$ of an FRW comoving diamond. The identification ``the volume in
HPS $=$ the volume in LL'' requires a dimensional reduction / uplift bridge that
is not in either paper.
\end{obstruction}

\begin{obstruction}[Vardian OAQEC type-I limitation, inherited from O1]\label{obs:T6}
Vardian's OAQEC tensor decomposition $\HH = \HH_A \otimes \HH_{\bar A}$
(needed for $K|_\AA$, eqs.~21, 23, 27) fails in type II$_\infty$ (no density
matrix for the diamond algebra). Identical to obstruction O1 of the
Vardian--Hollands run; survives unchanged in the trio attempt.
\end{obstruction}

\begin{remark}[UOGH gap survives]\label{rem:T4-residual}
Even waving T1--T3 hypothetically, the Krylov--Diameter Thm~4 proof sketch
(\texttt{krylov\_diameter.tex}, lines~347--360) explicitly uses
$\lambda_L^{\mathrm{mod}} = 2\pi$ from CMPT eq.~(4.16)/(5.4), which is the
type-III$_1$ universal operator-growth hypothesis on chiral CFT. The Block~A
definition (\texttt{krylov\_diameter.tex}, Def.~2) lifts this to type II$_\infty$
via the conditional expectation $E:\NN \to \AA(D_O)_\FRW$, but the slope
$2\pi$ remains \emph{inherited}, not proved intrinsically on II$_\infty$. Open
Question~1 of \texttt{krylov\_diameter.tex} (lines~485--497) explicitly flags
this. Neither HPS, LL, nor Vardian addresses UOGH on II$_\infty$ KMS states.
\end{remark}

\section{Recommendation}\label{sec:rec}

\begin{enumerate}
\item \textbf{Do not} write up the trio chain as a Krylov--Diameter supplement.
\item \textbf{Add} a one-paragraph defensive-citation appendix to
\texttt{paper/krylov\_diameter/krylov\_diameter.tex} listing T1--T6 with citations
to Vardian, HPS, LL.
\item \textbf{Cite} LL favourably as a programmatic reference for ``volume-as-index in
algebraic holography''; the dS application of LL \S{}V is the closest published
analogue to what one would want for FRW.
\item \textbf{Cite} HPS as the strongest published ``Krylov-spread $=$ bulk volume''
statement to date.
\item \textbf{Open problem to add to \texttt{krylov\_diameter.tex} Section ``Open
problems''}: \emph{Is there a subfactor inclusion}
$\AA(D_O')_\FRW \subset \AA(D_O)_\FRW$ \emph{between nested FRW comoving diamonds
whose Kosaki/Jones index equals} $\exp\!\big(C\cdot\Rprop(\eta_c)\big)$?
\emph{An affirmative answer, combined with an extension of HPS to type II$_\infty$,
would close the Krylov--Diameter proof.}
\end{enumerate}

Theorem~4 of \texttt{krylov\_diameter.tex} remains \emph{conditional} on UOGH
transfer (Remark~2, Open Question~1). The trio attempt \emph{reinforces} rather
than weakens that conditional stance.

\paragraph{arXiv API verification (spot-check, /tmp/\{vardian,hps,ll\}\_api.xml).}
\begin{itemize}
\item \texttt{2602.02675v1} (Vardian): \texttt{totalResults}=1, ``Modular Krylov
Complexity\dots'', published 2026-02-02T19:00:24Z, hep-th, 411\,866-byte PDF.
\item \texttt{2412.17785v2} (HPS): \texttt{totalResults}=1, ``Krylov spread
complexity as holographic complexity beyond JT gravity'', published
2024-12-23T18:43:35Z, journal\_ref \emph{Phys.\ Rev.\ Lett.}\ \textbf{135}
(2025) 151602, 501\,832-byte PDF.
\item \texttt{2508.00056v1} (Leutheusser--Liu): \texttt{totalResults}=1,
``Volume as an index of a subalgebra'', published 2025-07-31T18:00:01Z,
hep-th + gr-qc + math-ph, 45 pp., 1\,083\,703-byte PDF.
\end{itemize}

\end{document}
